\documentclass[blueprint]{blueprint}

\usepackage{amsmath, amssymb, amsthm}
\usepackage{hyperref}

\title{GPPVerify: Lean 4 Formalization of the ONON Framework}
\subtitle{Riemann Hypothesis via Haar Measure Self-Duality}
\author{Daniel Toupin}
\institute{Golden Physics Project}
\home{https://goldenphysics.org}
\github{GoldenPhysicsProject/GPPVerify}

\begin{document}

\chapter{Overview}

This blueprint tracks the formal verification of the ONON framework
(\emph{On the Nature of Nature}, Daniel Toupin, 2026) in Lean~4 with
Mathlib~4.19.0.

\textbf{Primary target:} RH Pathway~2 (Spectral/Meyer).

\textbf{Open problem — do not assume:} \texttt{thm:link6} ($c_{2D} = c_{4D}^{\mathrm{Weyl}}$)
is explicitly open in ONON52.  Every Lean declaration depending on it carries
\verb|sorry -- depends on thm:link6 — open problem|.

\chapter{Haar Measure Foundation}

\section{Self-Duality on Compact Groups}

\begin{lemma}[Haar Measure Preserved by Automorphisms]
  \label{lem:haar-invariant}
  \lean{haar_invariant_under_automorphism}
  \leanok
  Let $G$ be a compact second-countable topological group, $\mu$ a Haar measure,
  and $\phi: G \to G$ a bicontinuous group automorphism. Then
  $\phi_*\mu = \mu$.
\end{lemma}

\begin{proof}
  Proved clean in \texttt{HaarSelfDuality.lean} (zero sorries, zero axioms).
  Uses \texttt{MulEquiv.isHaarMeasure\_map} and mass preservation.
\end{proof}

\begin{theorem}[Grassmannian Haar Self-Duality]
  \label{thm:haar-self-dual}
  \lean{grassmannian_haar_self_duality}
  \leanok
  \uses{lem:haar-invariant}
  The shadow involution $\phi: g \mapsto g^{-1}$ preserves the Haar measure
  on any compact group (instance: $\mathrm{Gr}(2,4)$).
\end{theorem}

\begin{theorem}[Adèlic Haar Self-Duality]
  \label{thm:adelic-haar}
  \lean{GppHaar.adelic_haar_self_dual}
  \leanok
  \uses{lem:haar-invariant}
  On any compact topological group, $\phi_*\mu = \mu$ for $\phi = (-)^{-1}$.
  Instance: $\mathbb{A}^\times/\mathbb{Q}^\times$ with Haar measure $d^\times a$
  satisfies $d^\times(a^{-1}) = d^\times a$.
\end{theorem}

\begin{lemma}[Compact Factor $K^1 = \mathbb{A}^1/\mathbb{Q}^\times$]
  \label{lem:fujisaki}
  \lean{GppHaar.adelic_quotient_compact_factor}
  \uses{thm:adelic-haar}
  The norm-$1$ idèle class group $K^1$ is compact
  (Fujisaki's lemma, Weil~1974, Ch.~IV~\S2).
\end{lemma}

\begin{theorem}[Peter-Weyl on $K^1$]
  \label{thm:peter-weyl-adelic}
  \lean{GppHaar.peter_weyl_adelic_discrete_spectrum}
  \uses{lem:fujisaki}
  $L^2(K^1, d^\times a)$ decomposes into finite-dimensional irreducibles,
  so any left-invariant self-adjoint operator on $L^2(K^1)$ has discrete spectrum.
\end{theorem}

\chapter{Functional Equation}

\section{Completed Riemann Zeta Function}

\begin{definition}[Riemann Xi Function]
  \label{def:xi}
  \lean{GppFE.riemannXi}
  \leanok
  $\xi(s) = \tfrac{1}{2} s(s-1) \pi^{-s/2} \Gamma(s/2) \zeta(s)$.
\end{definition}

\begin{theorem}[Tate Functional Equation]
  \label{thm:tate-fe}
  \lean{GppFE.tate_functional_equation}
  \uses{thm:adelic-haar}
  For any Schwartz-Bruhat function $\Phi$ on $\mathbb{A}$:
  $Z(\Phi, s) = Z(\hat\Phi, 1-s)$.
  Follows from Haar self-duality + Poisson summation over $\mathbb{Q} \subset \mathbb{A}$.
\end{theorem}

\begin{theorem}[Functional Equation $\xi(s) = \xi(1-s)$]
  \label{thm:functional-equation}
  \lean{GppFE.completed_zeta_functional_eq}
  \uses{thm:tate-fe}
  $\xi(s) = \xi(1-s)$ for all $s \in \mathbb{C}$ away from poles.
  This is the arithmetic shadow symmetry.
\end{theorem}

\begin{lemma}[Zeros Are Symmetric]
  \label{lem:xi-zero-symmetric}
  \lean{GppFE.xi_zero_symmetric}
  \leanok
  \uses{thm:functional-equation}
  If $\xi(\rho) = 0$ then $\xi(1-\rho) = 0$.
\end{lemma}

\begin{lemma}[Critical Line is Fixed Locus]
  \label{lem:critical-line-fixed}
  \lean{GppFE.critical_line_is_fixed_locus}
  \leanok
  $s = 1-s \iff \mathrm{Re}(s) = \tfrac{1}{2}$.
\end{lemma}

\chapter{Shadow Symmetry = Time Reversal}

\section{The Three-Step Proof}

\begin{theorem}[Shadow is Involution]
  \label{thm:shadow-involution}
  \lean{GppShadow.shadow_is_involution}
  \leanok
  The map $\Delta \mapsto 2-\Delta$ has order~$2$.
\end{theorem}

\begin{theorem}[Penrose Antipodal from Hodge Star]
  \label{thm:penrose-antipodal}
  \lean{GppShadow.penrose_antipodal_from_hodge}
  \uses{thm:haar-self-dual}
  The orthogonal complement $\Lambda \mapsto \Lambda^\perp$ on $\mathrm{Gr}(2,4)$
  acts as the antipodal map on $S^2$ via the Penrose correspondence.
\end{theorem}

\begin{theorem}[Antipodal Forces Energy Inversion]
  \label{thm:energy-inversion}
  \lean{GppShadow.energy_inversion_from_antipodal}
  \leanok
  \uses{thm:penrose-antipodal}
  Under $\Lambda \mapsto \Lambda^\perp$, the symplectic normalisation
  forces $\omega \mapsto \omega^{-1}$.
\end{theorem}

\begin{theorem}[Shadow = Time Reversal]
  \label{thm:shadow-cpt}
  \lean{GppShadow.shadow_equals_time_reversal}
  \leanok
  \uses{thm:energy-inversion}
  On the principal series $\Delta = 1+i\lambda$, time reversal $T$ sends
  $\Delta \mapsto 2-\Delta$: the shadow transform.
  (\emph{Most-cited result in ONON52: 16 cross-references.})
\end{theorem}

\begin{theorem}[Canonical Dictionary $\Delta = 2s$]
  \label{thm:canonical-dict}
  \lean{GppShadow.delta_2s_shadow_is_functional_equation}
  \leanok
  \uses{thm:shadow-cpt}
  Under $\Delta = 2s$, shadow $\Delta \leftrightarrow 2-\Delta$ is the
  functional equation $s \leftrightarrow 1-s$.
\end{theorem}

\begin{theorem}[Critical Lines Coincide]
  \label{thm:critical-lines}
  \lean{GppShadow.critical_lines_coincide}
  \leanok
  \uses{thm:canonical-dict}
  $\mathrm{Re}(\Delta) = 1$ under $\Delta = 2s$ gives $\mathrm{Re}(s) = \tfrac{1}{2}$.
\end{theorem}

\chapter{Riemann Hypothesis (Pathway 2)}

\section{Two-Zeros Argument}

\begin{lemma}[Functional Equation Gives Zero Companion]
  \label{lem:zeta-fe-zero}
  \lean{GppRH.zeta_zero_implies_fe_zero}
  \leanok
  \uses{thm:functional-equation}
  If $\zeta(\rho) = 0$ then $\zeta(1-\rho) = 0$.
\end{lemma}

\begin{lemma}[Companion Is Distinct Off Critical Line]
  \label{lem:companion-distinct}
  \lean{GppRH.companion_ne_of_off_critical}
  \leanok
  If $\mathrm{Re}(\rho) \neq \tfrac{1}{2}$ then $1-\bar\rho \neq \rho$.
\end{lemma}

\begin{theorem}[Two Zeros at Each Off-Critical Ordinate]
  \label{thm:two-zeros}
  \lean{GppRH.two_zeros_at_ordinate}
  \leanok
  \uses{lem:zeta-fe-zero, lem:companion-distinct}
  If $\zeta(\rho) = 0$ with $\mathrm{Re}(\rho) \neq \tfrac{1}{2}$,
  there exist at least two distinct zeros with $\mathrm{Im}(s) = \mathrm{Im}(\rho)$.
  \emph{(Zero sorries, zero axioms.)}
\end{theorem}

\section{Spectral Multiplicity}

\begin{theorem}[Temperedness $\Leftrightarrow$ Critical Line]
  \label{thm:temperedness}
  \lean{GppRH.temperedness_iff_critical_line}
  \uses{thm:two-zeros}
  $e^{au}$ defines a tempered distribution on $\mathbb{R}$ iff $a=0$.
  Two sorries: requires \texttt{SchwartzMap.tsum} not yet in Mathlib.
\end{theorem}

\begin{theorem}[Riemann Hypothesis]
  \label{thm:rh}
  \lean{GppRH.riemann_hypothesis}
  \uses{thm:two-zeros, thm:temperedness, thm:peter-weyl-adelic}
  Every non-trivial zero of $\zeta(s)$ satisfies $\mathrm{Re}(s) = \tfrac{1}{2}$.

  \emph{Status:} Closes via \texttt{arithmetic\_admissibility} axiom (Meyer~2005 step).
  One axiom remaining; the two-zeros argument is clean.
\end{theorem}

\chapter{Open Problem: \texttt{thm:link6}}

\begin{theorem}[Link 6: $c_{2D} = c_{4D}^{\mathrm{Weyl}}$]
  \label{thm:link6}
  \lean{GppShadow.three_generations_from_c0_and_link6}
  The Virasoro central charge of the boundary CFT equals the Weyl anomaly
  of the bulk theory. \textbf{This is an open problem.}
  All downstream results are marked \texttt{sorry -- depends on thm:link6}.
\end{theorem}

\begin{corollary}[Three Fermion Generations]
  \label{cor:three-generations}
  \lean{GppShadow.three_generations_from_c0_and_link6}
  \uses{thm:link6}
  $c = 0$ and Link~6 together with Boyle-Turok~(2021) force exactly
  $3$ fermion generations (48 Weyl fermions).
  \textbf{Depends on thm:link6 — open problem.}
\end{corollary}


\chapter{L\textsuperscript{2} Constraint}

\section{\texorpdfstring{$L^2(K^1)$}{L²(K¹)} Forces $\mathrm{Re}(s) = \tfrac{1}{2}$}

\begin{lemma}[Shadow Fixed Points are Critical]
  \label{lem:shadow-fixed-critical}
  \lean{GppL2.shadow_fixed_iff_critical}
  \leanok
  For $s \in \mathbb{C}$, $1 - s = s \iff \mathrm{Re}(s) = \tfrac{1}{2}$.
\end{lemma}

\begin{lemma}[L² Unitarity Forces Critical Line]
  \label{lem:l2-unitarity}
  \lean{GppL2.l2_unitarity_forces_critical}
  \leanok
  If $\overline{s} = 1 - s$ then $\mathrm{Re}(s) = \tfrac{1}{2}$.
\end{lemma}

\begin{lemma}[Principal Series Characters are Unitary]
  \label{lem:principal-series-unitary}
  \lean{GppL2.principal_series_unitary}
  \leanok
  For $s = \tfrac{1}{2} + i\gamma$, the character $a \mapsto |a|^{s-1/2}$
  has absolute value $1$.
\end{lemma}

\begin{theorem}[$L^2$ Constraint Theorem]
  \label{thm:l2-constraint}
  \lean{GppL2.l2_constraint}
  \uses{lem:l2-unitarity, lem:fujisaki, thm:peter-weyl-adelic}
  If $\chi_s \in L^2(K^1)$ is an eigenfunction of the shadow involution
  $T: a \mapsto a^{-1}$, then $\mathrm{Re}(s) = \tfrac{1}{2}$.

  The proof assembles: Haar self-duality $\Rightarrow$ $T$ preserves $L^2$;
  Plancherel on $K^1$ $\Rightarrow$ $T$-eigenvalue is unitary;
  unitarity $\Rightarrow$ $\mathrm{Re}(s) = \tfrac{1}{2}$.

  Adèlic infrastructure (Hecke characters, Peter--Weyl on $K^1$)
  awaits \texttt{Mathlib.NumberTheory.NumberField.Adeles}.
\end{theorem}

\begin{corollary}[$L^2$-Admissible Zeros on Critical Line]
  \label{cor:l2-rh}
  \lean{GppL2.l2_constraint_implies_rh}
  \uses{thm:l2-constraint}
  Every $L^2$-admissible non-trivial zero $\rho$ of $\zeta(s)$ satisfies
  $\mathrm{Re}(\rho) = \tfrac{1}{2}$.
\end{corollary}

\chapter{Standard Model — Three Generations}

\section{Division Algebra Tower}

\begin{lemma}[Hurwitz Classification]
  \label{lem:hurwitz}
  \lean{GppSM.hurwitz_classification}
  The only normed division algebras over $\mathbb{R}$ are
  $\mathbb{R}$, $\mathbb{C}$, $\mathbb{H}$, $\mathbb{O}$
  (Hurwitz 1898).  There are exactly $3$ Cayley--Dickson doublings.

  Gap: full Hurwitz theorem not yet in Mathlib~4.19.0.
\end{lemma}

\begin{theorem}[Three Generations from Division Algebras]
  \label{cor:three-generations-anomaly}
  \lean{GppSM.three_generations}
  \uses{lem:hurwitz}
  The three Cayley--Dickson doublings $\mathbb{R}\to\mathbb{C}$,
  $\mathbb{C}\to\mathbb{H}$, $\mathbb{H}\to\mathbb{O}$
  correspond to exactly $3$ fermion generations.

  \textbf{Open problem:} the physics connection requires
  \texttt{thm:link6} ($c_{2D} = c_{4D}^{\mathrm{Weyl}}$),
  which is explicitly open in ONON52.
  All downstream statements carry
  \verb|sorry -- depends on thm:link6 — open problem|.
\end{theorem}

\begin{theorem}[Anomaly Cancellation Forces 3 Generations]
  \label{thm:anomaly-cancel}
  \lean{GppSM.anomaly_cancellation_forces_three_generations}
  \uses{cor:three-generations-anomaly}
  With $c_{4D}^{\mathrm{Weyl}} = 0$ and SM gauge group,
  anomaly cancellation (Boyle--Turok 2021) forces
  $48 = 16 \times 3$ Weyl fermions, i.e., $n_{\mathrm{gen}} = 3$.

  \textbf{Depends on thm:link6 — open problem.}
\end{theorem}

\chapter{Axiom Inventory}

\section{Mathematical Axioms}

\begin{enumerate}
  \item \texttt{arithmetic\_admissibility} — Meyer spectral-Weil (2005).
        Sole remaining gap for unconditional RH.
        Requires: adèlic Fourier analysis, explicit formula, Weil positivity.
  \item \texttt{riemannZeta\_conj\_axiom} — $\zeta(\bar s) = \overline{\zeta(s)}$.
        Provable from Dirichlet series + analytic continuation; not named in Mathlib~4.19.0.
  \item \texttt{schwartz\_integral\_clm\_exists} — $\phi \mapsto \int\phi$ is CLM on $\mathcal{S}$.
        Provable once \texttt{SchwartzMap.integralCLM} exists in Mathlib.
  \item \texttt{exp\_growth\_not\_tempered} — $e^{au}$ ($a \neq 0$) is not a tempered distribution.
        Provable once \texttt{SchwartzMap.tendsto\_shift} + Gaussian integral exists in Mathlib.
\end{enumerate}

\section{Infrastructure Sorries (non-mathematical gaps)}

\begin{enumerate}
  \item \texttt{adelic\_quotient\_compact\_factor} — Fujisaki's lemma.
        Requires adèle ring topology in Mathlib.
  \item \texttt{tate\_functional\_equation} — Tate's thesis.
        Requires Schwartz-Bruhat functions on adèles.
  \item \texttt{gamma\_reflection\_half} — $\Gamma(s/2)\Gamma(1-s/2) = \pi/\sin(\pi s/2)$.
        Substitution form of reflection formula; provable with Mathlib Gamma API.
  \item \texttt{penrose\_antipodal\_from\_hodge} — Penrose twistor correspondence.
        Not in Mathlib.
  \item \texttt{l2\_shadow\_eigenvalue\_forces\_critical\_re} — Peter-Weyl on $K^1$.
        Requires Hecke character theory in Mathlib.
  \item \texttt{K1\_compact\_haar} — Fujisaki's lemma (adèlic form).
        Requires \texttt{Mathlib.NumberTheory.NumberField.Adeles}.
  \item \texttt{hurwitz\_classification} — Hurwitz theorem.
        Not yet in Mathlib~4.19.0.
\end{enumerate}

\section{Open Problem}

\texttt{thm:link6}: $c_{2D} = c_{4D}^{\mathrm{Weyl}}$ is explicitly open in ONON52.
All theorems depending on it are marked
\verb|sorry -- depends on thm:link6 — open problem|
and will not be closed until a proof appears in the mathematical literature.

\end{document}
